
\section{Completing the Foundation--Concise bridge}
\label{sec:bridge}

\subsection{What is already stronger than the public summary suggests}

The current reverse translator contains a theorem of the following strength:
for every Concise derivation
\[
        d:\Gamma\vdash_{\mathsf{Ax}}\varphi
\]
and every assignment \(v\) of the \(n\) free Concise variables to Foundation
terms,
\[
   \Gamma^{v}\vdash_{\Lambda}\varphi^{v}.
\]
It handles arbitrary finite contexts, implication, both quantifiers, equality
substitution, and environments under binders.  A second theorem universally
closes the image when the Concise context is empty.  Thus the
Concise-to-Foundation \emph{open} transport is essentially implemented; what
is missing is a public wrapper specialized to \(\PRAAx\), the standard theory,
and the native-order comparison.

The syntax bridge also proves:
\[
\begin{aligned}
 \mathsf{toConciseCode}(\mathsf{toFoundationCode}(f))&=f,\\
 \mathsf{toConciseTerm}(\mathsf{toFoundationTerm}(t))&=t,\\
 \mathsf{toConciseFormula}(\mathsf{toFoundation}(\varphi))&=\varphi.
\end{aligned}
\]
Hence the Concise--Foundation--Concise composite is exact.

\subsection{Open sequents by closure and internalization}

The opposite direction is publicly exposed mainly for closed theory
theorems.  It can be extended without a new proof calculus.

Let
\[
   \Gamma=[\gamma_1,\ldots,\gamma_m]
\]
be a finite context of Foundation formulas with \(n\) free variables, and let
\(\phi\) be the conclusion.  Define the implication chain
\[
 \mathsf{chain}(\Gamma,\phi)
   =\gamma_1\to(\gamma_2\to\cdots(\gamma_m\to\phi)\cdots)
\]
and let \(\mathsf{close}_n\) universally quantify its free variables.

\begin{lemma}[Finite-context internalization]
A Foundation finite-context derivation
\[
   \Gamma\vdash_T\phi
\]
yields a theory theorem
\[
   T\vdash\mathsf{close}_n(\mathsf{chain}(\Gamma,\phi)).
\]
Conversely, opening the universal closure and applying the assumptions recovers
the finite-context derivation.
\end{lemma}

\begin{proof}
Repeatedly apply the Foundation deduction theorem to discharge the finite
context.  Then apply generalization to the free-variable family in the fixed
canonical order.  Conversely, specialize or open the quantifiers at the
distinct free variables and apply modus ponens once for each context formula.
All operations already exist in the Foundation finite-context and close-family
libraries.
\end{proof}

Transport the closed sentence with the current Foundation-to-Concise theorem.
A syntax lemma identifies its translation with the universal closure of the
Concise translation of the implication chain.  Repeated \code{Prf.allOpen}
opens the binders.  Repeated \code{Prf.impE} with hypotheses from the translated
context yields the conclusion.

\begin{theorem}[Open Foundation-to-Concise transport, reduced form]
Assume the translated images of the Foundation theory axioms have
\(\PRAAx\)-derivations and assume the source formula family is identified with
its code-normalized image.  Then
\[
   \Gamma\vdash_T\phi
   \quad\Longrightarrow\quad
   \Gamma^{\mathsf C}\vdash_{\PRAAx}\phi^{\mathsf C}
\]
for every finite open context and every fixed assignment of free variables to
distinct Concise variables.
\end{theorem}

\begin{proof}
Internalize and close the source derivation, transport the closed theorem,
open its quantifiers, and eliminate the implication chain.  The only
nonlogical work is the translation/closure commutation lemma and the
code-normalization identification.  Binder order is fixed once, so no choice
of fresh variables is involved.
\end{proof}

An alternative implementation generalizes the existing
\code{hilbertToConcise} induction to carry a target context.  Every case is
unchanged except that hypotheses are available by \code{Prf.hyp} and the
context is shifted under generalization.  The closure method is recommended
first because it reuses more existing code and gives a small theorem surface.

\subsection{The remaining code-normalization theorem}

For a Foundation primitive-recursive construction code \(c\), define
\[
       \mathsf{norm}(c)
        =\mathsf{toFoundationCode}(\mathsf{toConciseCode}(c)).
\]
This need not be syntactically equal to \(c\).  In particular, a positive-
arity zero code becomes a composition of the nullary zero code with an empty
tuple.  What is required is extensional equality.

\begin{theorem}[Foundation code normalization]
For every arity \(k\), code \(c\), and tuple of terms \(\vec t\), the
Foundation reference theory proves
\[
       c(\vec t)=\mathsf{norm}(c)(\vec t).
\]
The proof is uniform in \(\vec t\).
\end{theorem}

\begin{proof}
Use structural induction on \(c\).

\paragraph{Zero.}
For the nullary code the statement is reflexivity.  For positive arity,
\(\mathsf{norm}(c)\) is the composition of nullary zero with an empty argument
tuple.  The composition axiom reduces it to zero, as does the source zero code.

\paragraph{Successor and projection.}
Normalization is syntactically exact.

\paragraph{Composition.}
Use the defining equation for composition on both sides.  Apply the induction
hypothesis to the outer function and to every member of the finite inner tuple,
then use equality congruence.

\paragraph{Primitive recursion.}
Let the recursion variable be \(y\) and the remaining parameters be
\(\vec x\).  Induct on \(y\) with the quantifier-free equality
\[
 c(y,\vec x)=\mathsf{norm}(c)(y,\vec x).
\]
At zero, unfold both recursion codes and apply the induction hypothesis for the
base code.  At successor, unfold both step equations; use the arithmetic
induction hypothesis for the recursive values, the structural induction
hypothesis for the step code, and congruence for all parameters.  The invariant
is one equality and is therefore admissible in iPRA/Foundation PRA.
\end{proof}

This is the precise missing theorem identified by the unification plan.  It is
not a semantic appeal to standard natural numbers; it is an object-theory
derivation for every construction code.

\subsection{Lifting normalization to terms and formulas}

\begin{corollary}[Term normalization]
For every Foundation term \(t\), the reference theory proves
\[
  t=\mathsf{toFoundationTerm}(\mathsf{toConciseTerm}(t))
\]
under any assignment of its free and bound variables.
\end{corollary}

\begin{proof}
Induct on the term.  Variables are reflexive.  At a function application, use
code normalization and the induction hypotheses for all arguments.
\end{proof}

\begin{corollary}[Formula round trip]
For every equality-only Foundation formula \(\phi\), the theory proves
\[
 \phi\leftrightarrow
 \mathsf{toFoundation}(\mathsf{toConciseFormula}(\phi)).
\]
The result commutes with substitution, free-variable environments, and
quantifiers.
\end{corollary}

\begin{proof}
Induct on the formula.  Equality atoms use term normalization twice and
substitutivity of equality.  Falsity is reflexive.  Connectives use
intuitionistic congruence.  Under a quantifier, apply the induction hypothesis
to the opened body and use the existing all/ex congruence lemmas.
\end{proof}

For the native standard language with order, compose this equality-only result
with the already implemented theorem that replacing \(s<t\) by its comparison-
bit equation and re-embedding preserves each sentence up to a provable
biconditional.

\subsection{The desired public theorem}

After the preceding normalization lemmas, the bridge should expose:

\[
\boxed{
 \Gamma\vdash_{\mathsf{standardTheory}}\phi
 \quad\Longleftrightarrow\quad
 \Gamma^{\mathsf C}\vdash_{\PRAAx}\phi^{\mathsf C}
}
\]
for finite open contexts, with a documented variable environment.  The closed
equivalence becomes the \(n=0,\Gamma=[]\) specialization.

A Lean API should separate:
\begin{itemize}
\item equality-only reference formulas;
\item native formulas with order;
\item arbitrary assignments of free variables;
\item the canonical distinct-variable assignment;
\item closed sentences.
\end{itemize}
This avoids forcing downstream developments to unfold environment machinery.

\subsection{Why the bridge matters for PCFM}

PCFM should use \code{Prf PRAAx} as its canonical endpoint.  The completed
bridge then permits:
\begin{itemize}
\item importing Foundation arithmetization theorems as module leaves;
\item exporting PCFM endpoints to exact sentence encodings used by Busy Beaver
and consistency projects;
\item comparing arithmetic symbols once instead of maintaining duplicate
theorem libraries;
\item migrating one consumer at a time without a global backend rewrite.
\end{itemize}

No theorem packet should choose a backend merely because one arithmetic lemma
is missing.  The common module interface and bridge make backend choice an
implementation detail at the boundary.
